\lysh{36683}{4}

\begin{alist}
\item Αντικαθιστούμε στον τύπο της $f$, και συγκεκριμένα στον κλάδο $|x| + 2$ όπου $x = 0$ και βρίσκουμε: $f(0) = |0| + 2 = 2$.
Άρα η $C_f$ τέμνει τον άξονα $y^\prime y$ στο σημείο $M(0, 2)$.

\item
\begin{rlist}
\item Για $x = -1$ είναι: $f(-1) = |-1| + 2 = 1 + 2 = 3$.
Για $x = -2$ είναι: $f(-2) = |-2| + 2 = 2 + 2 = 4$.
Άρα η ημιευθεία $y = -x + 2$ διέρχεται από τα σημεία $A(-1, 3)$ και $B(-2, 4)$.
Για $x = 1$ είναι: $f(1) = |1| + 2 = 1 + 2 = 3$.
Για $x = 2$ είναι: $f(2) = |2| + 2 = 2 + 2 = 4$.
Άρα η ημιευθεία $y = x + 2$ διέρχεται από τα σημεία $\Gamma(1, 3)$ και $\Delta(2, 4)$.

Από τη γραφική παράσταση της $f(x)=|x|+2$ (η οποία αποτελείται από τις δύο ημιευθείες), διαπιστώνουμε ότι τα σημεία τομής της $C_f$ με την ευθεία $y = 3$ είναι τα $A(-1, 3)$ και $\Gamma(1, 3)$.

\item Τα σημεία $A(-1, 3)$ και $\Gamma(1, 3)$ έχουν αντίθετες τετμημένες και ίσες τεταγμένες. Άρα είναι συμμετρικά ως προς τον άξονα $y^\prime y$.
\end{rlist}

\item
\begin{rlist}
\item Η ευθεία $y = \alpha$ είναι μια ευθεία παράλληλη στον άξονα $x^\prime x$ και διέρχεται από το σημείο $(0, \alpha)$. Όπως διαπιστώνουμε (αφού το ελάχιστο της $f$ είναι το $2$), η ευθεία $y = \alpha$ τέμνει τη $C_f$ σε δύο σημεία αν και μόνο αν $\alpha > 2$.

\item Ο τύπος της $f$ γράφεται: $f(x) = |x| + 2, x \in \mathbb{R}$.
Οι τετμημένες των σημείων τομής της $C_f$ με την ευθεία $y = \alpha$ είναι οι λύσεις της εξίσωσης:
\[f(x) = \alpha \Leftrightarrow |x| + 2 = \alpha \Leftrightarrow |x| = \alpha - 2.\]

Αν $\alpha < 2 \Leftrightarrow \alpha - 2 < 0$, η εξίσωση $|x| = \alpha - 2$ είναι αδύνατη και επομένως η $C_f$ με την ευθεία $y = \alpha$ δεν έχουν κοινά σημεία.

Αν $\alpha = 2$, η εξίσωση $|x| = 2 - 2 \Leftrightarrow |x| = 0 \Leftrightarrow x = 0$ και επομένως η $C_f$ με την ευθεία $y = 2$ έχουν ένα κοινό σημείο το $(0, 2)$.

Αν $\alpha > 2$, η εξίσωση $|x| = \alpha - 2 \Leftrightarrow x = \alpha - 2$ ή $x = -(\alpha - 2) = -\alpha + 2$, δηλαδή δύο λύσεις διαφορετικές και επομένως η $C_f$ με την ευθεία $y = \alpha$ έχουν δύο κοινά σημεία, τα $(\alpha - 2, \alpha)$ και $(-\alpha + 2, \alpha)$.
\end{rlist}
\end{alist}
